% =====================================================================
%  1. Contextualização e objetivos
% =====================================================================
\section{Contextualização e objetivos}
\label{sec:contexto}

\subsection{Justificativa}

O Subcomponente 3.2 do PSH-PR tem como objetivo ampliar o acesso a serviços
sustentáveis de abastecimento de água e esgotamento sanitário no meio rural
paranaense, contribuindo para a universalização do saneamento, a promoção da
saúde pública e a proteção dos recursos hídricos \cite{mop2026}. A
universalização, estabelecida pelo novo marco legal do saneamento básico
\cite{lei14026}, ainda está distante nas áreas rurais do Estado: segundo dados
do \citeonline{sinisa2024}, a cobertura por rede de abastecimento de água
alcança 99,47\% da população urbana, mas apenas 15,11\% da população rural; no
esgotamento sanitário, a cobertura por rede coletora é de 84,41\% na área
urbana e de 2,04\% na área rural.

O saneamento rural parte do propósito de universalizar o acesso aos serviços
por meio de ações que garantam equidade, integralidade, intersetorialidade,
sustentabilidade dos serviços, participação e controle social. Diferentemente
do meio urbano, a alta dispersão populacional exige soluções técnicas --
coletivas ou individuais -- adaptadas à realidade local e acompanhadas de
modelos de gestão capazes de manter os sistemas em funcionamento após a sua
implantação \cite{mop2026}.

\pendencia{No MOP, o primeiro parágrafo do item B.3.2.2 menciona apenas
``universalizar o acesso ao Esgotamento Sanitário'', embora a ação trate
também de abastecimento de água. Há ainda um erro de digitação:
``soluções técnicas 0(coletivas ou individuais)''.}

\subsection{Objetivo geral}

Ampliar o acesso da população rural do Paraná ao abastecimento de água e ao
esgotamento sanitário por meio da implantação e da reabilitação de sistemas
coletivos de abastecimento de água, da instalação de sistemas descentralizados
de esgotamento sanitário e da implantação de modelos de gestão que assegurem
a sustentabilidade dos serviços.

\subsection{Objetivos específicos}

\begin{enumerate}[label=\alph*)]
  \item desenvolver modelos de gestão que permitam alcançar a
        universalização do saneamento rural com qualidade e sustentabilidade
        na prestação dos serviços públicos;
  \item implantar novos sistemas coletivos de abastecimento de água e
        sistemas descentralizados de esgotamento sanitário em comunidades
        rurais do Estado, visando ao princípio da integralidade;
  \item instalar sistemas descentralizados de esgotamento sanitário nas
        áreas prioritárias do PSH-PR de atuação do IDR-Paraná,
        preferencialmente nas áreas de mananciais de abastecimento público;
  \item reabilitar sistemas coletivos de abastecimento de água e substituir
        sistemas descentralizados de esgotamento sanitário em comunidades
        rurais, com foco na recuperação e na sustentabilidade.
\end{enumerate}

\subsection{Metas}
\label{subsec:metas}

Após a aplicação dos critérios de seleção definidos em âmbito estadual
(\autoref{sec:area}), a ação será conduzida em três frentes de trabalho
\cite{mop2026}:

\begin{enumerate}[label=\roman*)]
  \item implantar ou reabilitar sistemas de abastecimento de água em
        83 comunidades rurais (5.810 domicílios) e instalar sistemas
        descentralizados de esgotamento sanitário em 70\% dos domicílios
        atendidos (4.067 domicílios);
  \item aplicar modelos de gestão aos sistemas de abastecimento de água e de
        esgotamento sanitário implantados nas comunidades rurais e aos
        sistemas de abastecimento implantados pelo Programa Sanepar Rural;
  \item instalar 2.300 sistemas descentralizados de esgotamento sanitário,
        preferencialmente em mananciais de abastecimento público
        selecionados por critério qualitativo, na área prioritária do PSH-PR
        de atuação do IDR-Paraná.
\end{enumerate}

No conjunto, estima-se atender 22.708 habitantes com acesso à água potável e
ao esgotamento sanitário (\autoref{qua:indicadores}).

\pendencia{O MOP associa os 22.708 habitantes às 83 comunidades. Pelo fator
de 2,8 habitantes por domicílio (Censo 2022), adotado nos indicadores do
Componente~3, as 83 comunidades (5.810 domicílios) somam cerca de 16.268
habitantes. O total de 22.708 corresponde a (5.810 + 2.300) domicílios
$\times$ 2,8, ou seja, inclui os domicílios da área do IDR-Paraná. Ajustar o
texto do MOP.}
